% 03-wigner.tex -- rotation.  Formulas F4-F8.
% Input by mostlyharmless.tex.  No preamble here.

\section{Rotation}
\label{sec:wigner}

Changing the viewing geometry is a linear transformation of the coefficients, not
a refitting of the map.  This section states the law, the matrices that implement
it, and the non-negative-order form in which the transform is actually applied.
It is the one place where negative orders are wanted, and they appear only as
indices of the representation, never as coefficients.

\subsection{The rotation law}
\label{sec:rotation-law}

With $\hat n$ a direction in the body frame and $R$ the body-to-observer rotation,
the harmonics transform as
\begin{equation}
  Y_{lm}(R^{-1}\hat n) = \sum_{m'=-l}^{l} D^l_{m'm}(R) \, Y_{lm'}(\hat n) ,
  \label{eq:rotation-law}
\end{equation}
so a map whose coefficients are $A_{lm}$ in the body frame has coefficients
\begin{equation}
  a_{lm'} = \sum_{m=-l}^{l} D^l_{m'm}(R) \, A_{lm}
  \label{eq:coefficient-rotation}
\end{equation}
in the observer frame: the matrix contracts the \emph{second} index of $D$, and no
conjugation is involved.  The two symbols are kept apart throughout: capital $A_{lm}$ is
body-frame, lower case $a_{lm}$ is observer-frame, and the visibility of
sections~\ref{sec:visibility}--\ref{sec:full-map} always takes the latter, since that is
the frame the modes $V_{lm}$ are defined in.  Equation \eqref{eq:rotation-law} says that the two
descriptions are one physical map evaluated at rotated arguments --- a statement
that certif\-icate \vtag{6} tests directly.

Only non-negative orders need be carried, because the negative-order coefficients are
determined.  Splitting the sum in \eqref{eq:coefficient-rotation} by sign, and substituting
the conjugate for each $m < 0$, gives an expression in which every coefficient carries a
non-negative order:
\begin{equation}
  a_{lm'} = D^l_{m'0} \, A_{l0}
  + \sum_{m=1}^{l} \Big[ \, D^l_{m'm} \, A_{lm}
  + (-1)^m \, D^l_{m',-m} \, \overline{A_{lm}} \, \Big] ,
  \qquad m' \ge 0 .
  \label{eq:rotation-nonnegative}
\end{equation}
The $m = 0$ term appears once and not twice, because $A_{l0}$ is real: written as a pair it
would be doubled.  Every negative order in \eqref{eq:rotation-nonnegative} is an index of
$D$; none is a coefficient.  This is the one place the negative orders are needed as
\emph{indices}, and section~\ref{sec:parity} shows that the visibility needs no more than
their consequence.

One corollary is used by the deformation, so it is worth writing out here.  A zonal harmonic
about an axis $\hat e$ is $Y_{l0}$ in the frame whose pole is $\hat e$, and rotating it back
with \eqref{eq:coefficient-rotation} gives
\begin{equation}
  P_l(\hat n\cdot\hat e) \;=\; \frac{1}{2l+1}
    \sum_{m=-l}^{l} D^l_{m0}(R_{\hat e})\,Y_{lm}(\hat n) ,
  \label{eq:zonal-in-d}
\end{equation}
which is the addition theorem of the literature, obtained here without citing it.  Written
in stored orders, and at $l = 2$ expanded, with $\hat e = (a,b,c)$ and
$\hat n$ the observer direction,
\begin{equation}
  P_2(\hat n\cdot\hat e) \;=\; \frac{1}{5}
    \Bigl[\, \overline{Y_{20}(\hat e)}\,Y_{20}(\hat n)
    \;+\; 2\sum_{m=1}^{2} \mathrm{Re}\bigl(
    \overline{Y_{2m}(\hat e)}\,Y_{2m}(\hat n) \bigr) \Bigr] ,
  \label{eq:zonal-l2}
\end{equation}
and the five numbers it needs are
\begin{equation}
  \begin{aligned}
    Y_{20}(\hat e) &= \sqrt{\tfrac54}\,\bigl(3c^2 - 1\bigr) , \\
    Y_{2,\pm1}(\hat e) &= \mp\sqrt{\tfrac{15}{2}}\,c\,\bigl(a \pm ib\bigr) , \\
    Y_{2,\pm2}(\hat e) &= \sqrt{\tfrac{15}{8}}\,\bigl(a \pm ib\bigr)^2 .
  \end{aligned}
  \label{eq:harmonic-at-axis-l2}
\end{equation}
So a zonal quadrupole about a known axis contributes these five coefficients and no
summation over $m$ is needed anywhere in the deformation.  Certificate \vtag{3} checks
\eqref{eq:harmonic-at-axis-l2} against the general \code{sph_harm_addition}.

\subsection{The Wigner matrices}
\label{sec:wigner-matrices}

In the zyz convention, with $R = R_z(\alpha) R_y(\beta) R_z(\gamma)$,
\begin{equation}
  D^l_{m'm}(\alpha,\beta,\gamma) = e^{-im'\alpha} \, d^l_{m'm}(\beta) \, e^{-im\gamma} ,
  \qquad
  d^l_{m'm}(\beta) = \langle l m' | e^{-i\beta J_y} | l m \rangle ,
  \label{eq:wigner-D}
\end{equation}
the $d$ matrix is real, and its explicit form is
\begin{equation}
\begin{split}
  d^l_{m'm}(\beta) = {}&
  \sqrt{(l+m')!\,(l-m')!\,(l+m)!\,(l-m)!} \\
  &\times \sum_s \frac{(-1)^{m'-m+s}}{(l+m-s)!\,s!\,(m'-m+s)!\,(l-m'-s)!} \\
  &\times \cos^{2l+m-m'-2s}\!\tfrac{\beta}{2}\;
  \sin^{m'-m+2s}\!\tfrac{\beta}{2},
\end{split}
  \label{eq:wigner-d-explicit}
\end{equation}
summed over the $s$ for which every factorial is defined.  Certificate \vtag{4}
checks \eqref{eq:wigner-d-explicit} against the exponential in
\eqref{eq:wigner-D}, which is an independent construction of the same object.
\emph{Numerical}, $l \le 6$.

The elementary properties follow from \eqref{eq:wigner-D} and are collected in
certificate \vtag{5}: the $d$ matrix is real and orthogonal,
$d^l_{m'm}(0) = \delta_{m'm}$, $d^l_{00}(\beta) = P_l(\cos\beta)$, and
$d^l_{m'm} = (-1)^{m-m'} d^l_{mm'}$.  \emph{Numerical}, $l \le 6$.

Only the two diagonal factors in \eqref{eq:wigner-D} are complex, and they carry
the spin phase and the position angle; the inclination enters only through the
real $d$ matrix.

\subsection{The viewing angles}
\label{sec:viewing-angles}

The body-to-observer rotation is built from the spin phase $\alpha$, the
inclination $\beta$ and the position angle $\PA$ of the spin axis.  Because the
observer's $x$ axis is North (D24), $\PA$ is measured from North through East
directly, so no offset appears anywhere:
\begin{equation}
  R(\alpha,\beta,\PA) = R_z(\PA)\, R_y(\beta)\, R_z(\alpha) ,
  \label{eq:body-to-observer-wigner}
\end{equation}
which is \eqref{eq:wigner-D} with $\alpha_{\rm zyz} = \PA$ and
$\gamma_{\rm zyz} = \alpha$:
\begin{equation}
  D^l_{m'm}(R) = e^{-im'\PA} \, d^l_{m'm}(\beta) \, e^{-im\alpha} .
  \label{eq:D-viewing}
\end{equation}
Note which angle multiplies which index: $\PA$ goes with the row index $m'$, and
$\alpha$ with the column index $m$.  The middle axis stays $y$: the inclination is
applied first and the position angle then carries the result to its place on the
sky.  The body $z$ axis, which is the spin, maps to
$(\sin\beta\cos\PA,\ \sin\beta\sin\PA,\ \cos\beta)$, so it tilts toward the azimuth
$\PA$ and sits $\cos\beta$ above the sky plane.

\subsection{What was checked}
\label{sec:wigner-register}

Equation \eqref{eq:rotation-nonnegative} is checked against the full-range
transform \eqref{eq:coefficient-rotation} by \vtag{6}, and the full-range
transform against direct geometry by the same certificate: rotating the
coefficients must agree with evaluating the original map at the inversely rotated
direction,
\begin{equation}
  I_{\rm new}(\hat n) = I_{\rm body}(R^{\!\top}\hat n) ,
  \label{eq:rotation-geometry}
\end{equation}
which shares nothing with the algebra.  Composition is checked the same way, by
applying two rotations and comparing against the composite direction.
\emph{Numerical}, $l \le 4$ and $l \le 3$ respectively.

\begin{center}
\begin{tabular}{lll}
\hline
tag & claim & rung \\
\hline
\vtag{4} & $d$ from \eqref{eq:wigner-d-explicit} against $\exp(-i\beta J_y)$ & numerical, $l \le 6$ \\
\vtag{5} & properties of $d$ & numerical, $l \le 6$ \\
\vtag{6} & \eqref{eq:rotation-nonnegative} and \eqref{eq:rotation-geometry} & numerical, $l \le 4$ \\
\vtag{7} & composition of two rotations & numerical, $l \le 3$ \\
\hline
\end{tabular}
\end{center}
